










\subsubsection{In the purified and mixed sample access models: collective case}
\label{sec:fidelity estimation collective}



We now provide the square root fidelity estimation algorithms in the purified and mixed sample access models when collective operations over multiple samples of the states are allowed.
Our statement is as follows.



\begin{theorem}[Square root fidelity estimation in the purified and mixed sample access models with collective operations]
\label{thm:fidelity est local sample collective}
Let $\delta \in (0, 1)$. Then, in each of the purified and mixed sample access models, there is a quantum sample algorithm that outputs $\sqrt{\til{\rF}}$ such that $\big|\sqrt{\til{\rF}} - \sqrt{\rF}\big| \leq \delta$ with probability at least $2/3$. The algorithm uses $n = \cO\big(\f{1}{\delta^2}{\rm min}\big\{\f{1}{s_{\rm min}^2}, \f{r^2}{\delta^2}\big\}\big(\log{\big(\f{1}{\delta}\big)}\big)^2\big)$ samples of $\ket{\rho}^{\sA\sB}$ and $\ket{\sigma}^{\sA\sB}$ in the purified sample access model, or $\rho^\sA$ and $\sigma^\sA$ in the mixed sample access model.
The quantum circuit of this algorithm consists of $\til{\cO}(n^4)$ one- and two-qubit gates, and $\til{\cO}(n^2)$ qubits suffice at any one time.

\end{theorem}





The key idea in the construction of the algorithms in \Cref{thm:fidelity est local sample collective} is to use a technique recently described in Ref.~\cite{tang2025conjugatequerieshelp}, which aims to simulate a query algorithm by a sample algorithm. This is achieved by using sampled states to approximately implement unitaries that prepare their purified states.
If the original query algorithm attains its goal with probability at least $1-\eta$ using $q$ queries, the corresponding sample algorithm can simulate the same goal with probability at least $1-\eta-\epsilon$ using $\cO(q^2/\epsilon)$ samples.
With this technique, all $q = \cO\big(\f{1}{\delta}{\rm min}\big\{\f{1}{s_{\rm min}}, \f{r}{\delta}\big\}\log{\big(\f{1}{\delta}\big)}\big)$ queries to the purified state preparation unitaries $U_\rho^{\sA\sB}$, $U_\sigma^{\sA\sB}$, and their inverses in the query algorithm $\mathtt{UhlFidelityEstPurifQuery}$ (\Cref{thm:fid est quel}) are replaced by approximate operations constructed from multiple copies of $\ket{\rho}^{\sA\sB}$ and $\ket{\sigma}^{\sA\sB}$, or $\rho^\sA$ and $\sigma^\sA$. Consequently, we then obtain the result in \Cref{thm:fidelity est local sample collective}.


We formally introduce the technique in Ref.~\cite{tang2025conjugatequerieshelp}, slightly adapted for our purposes.
In the following, we focus on the discussion of the mixed sample access model.
For the purified sample access model, the step of random purification described below can simply be omitted.




\begin{theorem}[Simulating queries to state preparation unitary given copies of the state~{\cite[Theorem 1.5]{tang2025conjugatequerieshelp}}]
\label{thm:sample to query lift}
Let $\epsilon, \eta \in (0, 1)$ and let $U_\omega^{\sA\sB}$ be a unitary such that $U_\omega^{\sA\sB}\ket{0}^{\sA\sB} = \ket{\omega}^{\sA\sB}$. Suppose that there is a quantum query algorithm that makes $q$ queries to $U_\omega^{\sA\sB}$ and $(U_\omega^{\sA\sB})^\dag$, and produces an output with probability $p^{\rm qry}$.
If the query algorithm satisfies $p^{\rm qry} \geq 1 - \eta$ for any choice of the purifying system $\sB$, then there exists a quantum sample algorithm that simulates the output of the query algorithm with probability $p^{\rm samp} \geq 1 - \eta - \epsilon$, using $n = \cO(q^2 / \epsilon)$ samples of the state $\omega^\sA$.
The quantum circuit of the sample algorithm consists of $\til{\cO}(n^4)$ one- and two-qubit gates, and $\til{\cO}(n^2)$ qubits at any one time, in addition to those of the query algorithm.



\end{theorem}



At the heart of this algorithm are the following three primitives.
The first primitive is called random purification, which is regarded as an algorithmic strengthening of the results in Refs.~\cite{Soleimanifar2022testingmatrixproduct, chen2024localtestUniInv}. 
The term \emph{random} refers to the following: let $\ket{\omega_0}^{\sA\sB}$ be a fixed purified state of $\omega^\sA$. Then, all purified states of $\omega^\sA$ on $\sA\sB$ are connected by a unitary acting on the purifying system $\sB$; that is, for any purified state, there is a unitary $U^\sB$ such that it can be written as $U^\sB\ket{\omega_0}^{\sA\sB}$.
The random means that the unitary $U^\sB$ is drawn uniformly at random according to the Haar measure. 
The random purification subroutine then allows us to approximately map $(\omega^\sA)^{\otimes n}$ to a state $\bE_U[(U^\sB\ketbra{\omega_0}{\omega_0}^{\sA\sB}(U^{\sB})^\dag)^{\otimes n}]$, i.e., to a state uniformly averaged over the choice of purification.
Remarkably, although the state is averaged, the algorithm prepares $n$ purified states from $n$ samples of $\omega^\sA$ without reducing the number of samples.
The primary component of this subroutine is the Schur transform~\cite{harrow2005applicationSchurtrans, Bacon2006efficientcircuitSchur, Krovi2019efficienthigh, Nguyen2024theorynuralnet, burchardt2025highdimensionalschur}, which mainly determines the number of additional gates and qubits to the circuit of the query algorithm as $\til{\cO}(n^4)$ and $\til{\cO}(n^2)$~\cite{burchardt2025highdimensionalschur}, respectively.


The second primitive, which appeared earlier in Refs.~\cite{kretschmer2021pseudorandomnessclassical, Goldin2025translatingcommonhaar}, relates to a state with a phase: $\ket{\omega_\theta}^{\sA\sB\sS} = \f{1}{\sqrt{2}}(e^{i\theta}\ket{0}^{\sA\sB}\ket{0}^{\sS} + \ket{\omega}^{\sA\sB}\ket{1}^{\sS})$, where $\sS$ is a one-qubit auxiliary system. 
Using this primitive, we can map $(\ketbra{\omega}{\omega}^{\sA\sB})^{\otimes n}$ to $\bE_\theta[(\ketbra{\omega_\theta}{\omega_\theta}^{\sA\sB\sS})^{\otimes n}]$, where the expectation $\bE_\theta[\cdot]$ is taken over $\theta$ chosen uniformly at random from $[0, 2\pi)$. 
The quantum circuit implementing this primitive consists of $\cO(n^2\log{(d_\sA d_\sB)})$ one- and two-qubit gates and uses $\cO(n\log{(d_\sA d_\sB)})$ qubits at any one time \cite{tang2025conjugatequerieshelp}.


The last primitive is the density matrix exponentiation (\Cref{prevthm:cntrl-DME} with $t=\pi$), which maps $(\ketbra{\omega_\theta}{\omega_\theta}^{\sA\sB\sS})^{\otimes m}$ to $e^{i\pi \ketbra{\omega_\theta}{\omega_\theta}^{\sA\sB\sS}}$ within a diamond norm error $\epsilon'$, where $m = \cO(1/\epsilon')$. 
Note that the unitary $e^{i\pi \ketbra{\omega_\theta}{\omega_\theta}^{\sA\sB\sS}}$ is a purified state preparation unitary because $e^{i\pi \ketbra{\omega_\theta}{\omega_\theta}^{\sA\sB\sS}} \ket{0}^{\sA\sB\sS} = -e^{-i\theta} \ket{\omega}^{\sA\sB} \ket{1}^\sS$. Here, the state $-e^{-i\theta} \ket{\omega}^{\sA\sB} \ket{1}^\sS$ is one of the purified states of $\omega^\sA$ with purifying system $\sB\sS$ of dimension at least twice the rank of $\omega^\sA$.
To apply the unitary for $q$ times with total error at most $\epsilon$, we set $\epsilon' = \epsilon/q$ and use $n = qm = \cO(q^2/\epsilon)$ copies of $\ket{\omega_\theta}^{\sA\sB\sS}$.
Naturally, the inverse of the purified state preparation unitary is obtained in a similar manner.


Now, consider a query algorithm that outputs a certain result with probability $p^{\rm qry}$, using $q$ queries to a unitary preparing a purified state of $\omega^\sA$.
We suppose that this query algorithm is independent of the choice of purification, i.e., it satisfies $p^{\rm qry} \geq 1 - \eta$ for any given purified state preparation unitary.
Without loss of generality, we then consider a sufficiently large purifying system, at least twice the rank of $\omega^\sA$.
By replacing all $q$ queries in the query algorithm with the three primitives above, we obtain a sample algorithm.
Using $n = \cO(q^2/\epsilon)$ copies of the state $\omega^\sA$, the sample algorithm can simulate the original query algorithm with probability $p^{\rm samp}$, such that $\big|p^{\rm samp} - \bE_\theta\bE_U [p^{\rm qry}]\big| \leq \epsilon$.
Since $p^{\rm qry} \geq 1 - \eta$ for any $U^\sB$ and $\theta$, it holds that $p^{\rm samp} \geq \bE_\theta \bE_U [p^{\rm qry}] - \epsilon \ge 1 - \eta - \epsilon$.
Hence, \Cref{thm:sample to query lift} follows.





As mentioned in Ref.~\cite{tang2025conjugatequerieshelp}, we should note a caveat: the query algorithm to which we apply \Cref{thm:sample to query lift} must work on any choice of purified state preparation unitary. Since the first and second primitives involve the averages $\bE_\theta[\cdot]$ and $\bE_U[\cdot]$ over possible purifications, the result cannot depend on a specific purification. 
While this condition is typically satisfied in common query algorithms, depending on tasks, a particular purification may be more advantageous than others. For instance, this is the case for the Uhlmann transformation itself. Specifically, the Uhlmann transformation works only for the specific purifications determined by the given initial and final states, $\ket{\rho}^{\sA\sB}$ and $\ket{\sigma}^{\sA\sB}$, and not for arbitrary purifications. On the other hand, when the goal is to estimate the fidelity, this issue can be circumvented by applying the purified state preparation unitaries $U_\rho^{\sA\sB}$ and $U_\sigma^{\sA\sB}$ for the initial and final states, along with the Uhlmann transformation (see also Eqs.~\eqref{inteq:170} to~\eqref{inteq:48}). For this reason, the sample cost of fidelity estimation in the mixed sample access model can be improved by this approach, whereas that of the Uhlmann transformation itself cannot.





Given the above results, \Cref{thm:sample to query lift}, we can straightforwardly prove \Cref{thm:fidelity est local sample collective}.
Note that, unlike other algorithms discussed in this paper, the algorithm in \Cref{thm:sample to query lift}---and thus our algorithms in \Cref{thm:fidelity est local sample collective}---need to handle $n$ samples of the state collectively.







\begin{proof}[Proof of Theorem~\ref{thm:fidelity est local sample collective}]
    

Recalling that for $\delta \in (0, 1)$, and $\eta \in (0, 1)$ of constant order, the algorithm
\begin{align}
    \mathtt{UhlFidelityEstPurifQuery}(U_\rho^{\sA\sB}, U_\sigma^{\sA\sB}; \delta)
\end{align}
in \Cref{alg:fidelestquery} outputs an estimate $\sqrt{\til{\rF}}$ that satisfies $\big|\sqrt{\til{\rF}} - \sqrt{\rF}\big| \leq \delta$ with success probability $p^{\rm qry} \ge 1-\eta$, using $q = \cO\big(\f{1}{\delta}{\rm min}\big\{\f{1}{s_{\rm min}}, \f{r}{\delta}\big\}\log{\big(\f{1}{\delta}\big)}\big)$ queries to $U_\rho^{\sA\sB}$, $U_\sigma^{\sA\sB}$, and their inverses.
Although the success probability was set to at least $2/3$ in \Cref{thm:fid est quel}, it can be amplified to $1-\eta$ with only a constant increase of the query cost as long as $\eta$ is of constant order (see Sec.~\ref{sec:fidelity estimation PQA}).
Note that the query algorithm $\mathtt{UhlFidelityEstPurifQuery}$ works on \emph{any} choice of purification, and hence, outputs the desired estimate with probability at least $1 - \eta$ independent of which purification is selected.


Following Theorem~\ref{thm:sample to query lift}, we implement the query algorithm with $q$ queries by the corresponding sample algorithm by approximately constructing the purified state preparation unitaries.
Using $n = \cO(q^2/\epsilon)$ samples of the states $\rho^\sA$ and $\sigma^\sA$, this sample algorithm outputs $\sqrt{\til{\rF}}$ satisfying $\big|\sqrt{\til{\rF}} - \sqrt{\rF}\big| \leq \delta$ with probability $p^{\rm samp}$ such that $p^{\rm samp} \geq 1 - \eta - \epsilon$.
By setting $\eta = 1/6$ and $\epsilon = 1/6$, we conclude that, for any $\delta \in (0, 1)$, the sample algorithm outputs $\sqrt{\til{\rF}}$ that satisfies $\big|\sqrt{\til{\rF}} - \sqrt{\rF}\big| \leq \delta$ with probability $p^{\rm samp} \geq 2/3$, using $n$ samples of the states $\rho^\sA$ and $\sigma^\sA$, where
\begin{align}
    n 
    &= \cO(q^2/\epsilon) \\
    &=\Big(\f{1}{\delta^2}{\rm min}\Big\{\f{1}{s_{\rm min}^2}, \f{r^2}{\delta^2}\Big\}\Big(\log{\Big(\f{1}{\delta}\Big)}\Big)^2\Big).
\end{align}




We now evaluate the number of gates and qubits in the quantum circuit of this sample algorithm.
The original query algorithm uses $\cO\big(q\log(d_\sA d_\sB) + (\log(1/\delta))^2\big)$ gates and $\cO(\log(d_\sA d_\sB) + \log(1/\delta))$ qubits at once. (\Cref{thm:fid est quel}). 
In addition, $n\log(d_\sA d_\sB)$ qubits are required to store the $n$ samples, so the number of qubits becomes $n\log(d_\sA d_\sB) + \cO(\log(d_\sA d_\sB) + \log(1/\delta)) = \cO(n\log(d_\sA d_\sB))$.
From \Cref{thm:sample to query lift}, converting the query algorithm to the sample one takes $\til{\cO}(n^4)$ gates and $\til{\cO}(n^2)$ qubits simultaneously, in the part of the Schur transform~\cite{burchardt2025highdimensionalschur}.
Therefore, the total number of one- and two-qubit gates in the whole sample algorithm is determined by their sum: $\cO\big(q\log(d_\sA d_\sB) + (\log(1/\delta))^2\big) + \til{\cO}(n^4) = \til{\cO}(n^4)$.
The maximum number of qubits used at one time is given by $\max\{\cO(n\log(d_\sA d_\sB)), \til{\cO}(n^2)\} = \til{\cO}(n^2)$.





\end{proof}





